% from intro-fr.tex
\oldpage[I]{5}

Ce mémoire, et les nombreux autres qui doivent lui faire suite, sont destinés à former un traité sur les fondements de la Géométrie algébrique.
Ils ne présupposent en principe aucune connaissance particulière de cette discipline, et il s'est même avéré qu'une telle connaissance, malgré ses avantages évidents, pouvait parfois (par l'habitude trop exclusive du point de vue birationnel qu'elle implique) être nuisible à celui qui désire se familiariser avec le point de vue et les techniques exposés ici.
Par contre, nous supposerons que le lecteur a une bonne connaissance des sujets suivants~:

\begin{enumerate}
  \item[a)] L'\emph{Algèbre commutative}, telle qu'elle est exposée par exemple dans les volumes en cours de préparation des \emph{Éléments} de N.~Bourbaki (et, en attendant la parution de ces volumes, dans Samuel-Zariski~\cite{I-13} et Samuel~\cite{I-11},~\cite{I-12}).
  \item[b)] L'\emph{Algèbre homologique}, pour laquelle nous renvoyons à Cartan-Eilenberg~\cite{I-2} (cité~(M)) et Godement~\cite{I-4} (cité~(G)), ainsi qu'à l'article récent de A.~Grothendieck~\cite{I-6} (cité~(T)).
  \item[c)] La \emph{Théorie des faisceaux}, où nos principales références seront~(G) et~(T); cette dernière théorie fournit le langage indispensable pour interpréter en termes <<~géométriques~>> les notions essentielles de l'Algèbre commutative, et pour les <<~globaliser~>>.
  \item[d)] Enfin, il sera utile au lecteur d'avoir une certaine familiarité avec le \emph{langage fonctoriel}, qui sera constamment employé dans ce Traité, et pour lequel le lecteur pourra consulter~(M),~(G) et surtout~(T); les principes de ce langage et les principaux résultats de la théorie générale des foncteurs seront exposés plus en détail dans un ouvrage en cours de préparation par les auteurs de ce Traité.
\end{enumerate}

\bigskip
\begin{center}*\quad*\quad*\end{center}
\bigskip

Ce n'est pas le lieu, dans cette Introduction, de donner une description plus ou moins sommaire du point de vue des <<~schémas~>> en Géométrie algébrique, ni la longue liste des raisons qui ont rendu nécessaire son adoption, et en particulier l'acceptation systématique d'éléments nilpotents dans les anneaux locaux des <<~variétés~>> que nous considérons (ce qui, nécessairement, relègue au second plan la notion d'application rationnelle, au profit de celle d'application régulière ou <<~morphisme~>>).
Le présent Traité vise précisément à développer de façon systématique le langage des <<~schémas~>> et démontrera, nous l'espérons, sa nécessité.
Encore qu'il serait facile de le faire, nous

